% !TEX root = ../main.tex
\chapter{Methodology}
\label{ch:method}

% =====================================================================
% 骨架來源：docs/thesis/thesis-outline.md §5 —— 全篇最有研究味的一章
%
% 【本章與第 3 章的分界】
%   第 3 章 = 靜態存在的東西（規格長什麼樣、怎麼被執行）
%   第 4 章 = 它怎麼被產生出來的
% 對應樣板原本的 "System Model" -> "<你的演算法>" 結構，本章就是那個演算法。
% =====================================================================

\section{State Machine}
\label{sec:der:statemachine}

% TODO derive -> gate -> oracle -> repair -> converge / escalate to human

\section{Structural Gates}
\label{sec:der:gates}

% TODO schema 合法性、locator 可解析、覆蓋率門檻、內部一致性。
%      重點論述：這些在花錢跑實例之前就能擋掉大部分壞規格。
% TODO 【已量到的證據】一份壞規則不得摧毀產生它的 derive；
%      gate 擋下的東西要能被人審查而不是靜默丟棄。

\section{Two Complementary Oracles}
\label{sec:der:oracles}

% TODO ★ 本論文的核心實證洞見之一 ★
%      derive-time expected-text oracle：規格宣稱「這個欄位會產生什麼文字」，立即比對
%      case-level oracle：實際執行後與參考輸出比對
%      兩者抓到的錯誤類型不同，任何一方都無法取代另一方。
%      【要給證據與失敗案例分類，第 5 章有對應的消融實驗】
%
% TODO 為什麼不用 LLM-as-judge：在沉默失敗面前，judge 的錯誤與生成的錯誤相關。
%      這條論述同時支撐第 5 章「評測器不得詢問另一個 LLM」的決定。

\section{Case-Repair Loop}
\label{sec:der:repair}

% TODO 失敗案例如何回饋成修正指令；收斂性與停止條件；避免震盪（同一個 gate 反覆修）。
% TODO 反過擬合守則：修正僅限值層，不允許為了通過某個案例而改動結構層。

\section{Human in the Loop: Evidence Review}
\label{sec:der:evidence}

% TODO 人力從「每份文件都檢查」移轉為「一次性審查證據」。
%      要交代介面呈現什麼證據（欄位、locator 命中位置、預期值 vs 實際值）。

\section{Coverage Is the Real Bottleneck}
\label{sec:der:coverage}

% TODO ★ 核心實證洞見之二 ★
%      在達到高準確率之後，剩餘風險主要來自「規格沒涵蓋到的欄位」而不是「涵蓋了但算錯」。
%
% TODO 【現況與承諾的落差，必須誠實處理】
%      目前分母 100% 由 agent 自陳（rule_inventory），分子也是同一個 agent 的輸出。
%      後果實測得到：同一份實作跨兩輪 derive 報出 80% -> 83%「進步」，
%      原因只是第二輪把兩筆併成一筆。更根本的是：少枚舉一條規則分數就變高。
%
%      修法是 rule_segmenter（確定性切段，以內容雜湊為身分），讓規則文件成為權威分母。
%      實作未開始（TODO.md P0 FEATURE）。
%      若到寫作時仍未完成，本節必須把「規則文件是覆蓋率的權威來源」寫成規範性主張，
%      並在第 5 章 Threats to Validity 明列現行計分的偏誤方向。
%      【不可以照抄 deck 的說法】deck 承諾的是「可以對照解析後的規則文件實際計數」。
